% =============================================================================
% EXAMPLE RESEARCH PLAN / PROPOSAL TEMPLATE
% =============================================================================
%
% COMPILER: pdfLaTeX (recommended)
%
% REQUIRED PACKAGES (install via TeX Live or MiKTeX):
%   - geometry     : 1-inch margins (required by most grants)
%   - fontenc      : T1 font encoding
%   - bold-extra   : Bold small caps for section headers
%   - xcolor       : Custom colors (garnet branding)
%
% OPTIONAL BUT RECOMMENDED:
%   - microtype    : Improved typography (letter spacing, etc.)
%   - hyperref     : Clickable cross-references (add if needed)
%
% COMPILE COMMAND:
%   pdflatex example_plan.tex
%
% FORMATTING NOTES:
%   - Single-spaced (no setspace package needed)
%   - 12pt font minimum (grant requirement)
%   - 1-inch margins on all sides (grant requirement)
%   - Unlocked, searchable PDF (default with pdflatex)
%
% CUSTOM COMMANDS:
%   \sect{Title}    - Section header (garnet bold small caps with period)
%   \subsect{Title} - Subsection header (garnet italic with period)
%
% =============================================================================

\documentclass[12pt]{article}

% --- Page Layout (1-inch margins required by most grants) ---
\usepackage[margin=1in]{geometry}

% --- Font Encoding & Bold Small Caps ---
\usepackage[T1]{fontenc}
\usepackage{bold-extra}  % Required for bold small caps

% --- Colors (UofSC Brand) ---
\usepackage{xcolor}
\definecolor{garnet}{HTML}{73000A}

% --- Tables ---
% Use p{0.XX\textwidth} for manual column widths that sum to ~0.95-1.0\textwidth

% --- Paragraph Formatting ---
\setlength{\parindent}{0.25in}  % Half-tab indent
\setlength{\parskip}{0.5em}     % Compact paragraph spacing

% --- Custom Section Commands ---
% These automatically handle \noindent to prevent section header indentation
\newcommand{\sect}[1]{\noindent{\color{garnet}\textbf{\textsc{#1}}.}}
\newcommand{\subsect}[1]{\noindent{\color{garnet}\textit{#1.}}}

% =============================================================================
\begin{document}

% --- Title Block ---
% Garnet bold small caps with full-width underline
\noindent{\color{garnet}\large\textbf{\textsc{Development of Autonomous Multi-Agent Coordination Framework for Distributed Sensor Networks}}}\\[-0.8em]
{\color{garnet}\rule{\textwidth}{0.5pt}}\\[0.2em]
\sect{Project Overview} Modern distributed sensor networks face significant challenges in coordination, fault tolerance, and real-time decision making. As sensor deployments scale from hundreds to thousands of nodes, centralized control architectures become bottlenecks for latency-critical applications. This project proposes a decentralized multi-agent coordination framework that enables autonomous decision-making at the edge while maintaining global coherence.

The proposed framework addresses three fundamental challenges in distributed sensing: (1) achieving consensus among agents with partial observability, (2) adapting to dynamic network topology changes, and (3) optimizing resource allocation under communication constraints. By combining techniques from distributed systems, game theory, and reinforcement learning, this work aims to establish theoretical foundations and practical algorithms for next-generation sensor networks.

Preliminary simulations demonstrate that decentralized approaches can reduce communication overhead by 60\% compared to centralized alternatives while maintaining 95\% decision accuracy. The proposed research extends these results to heterogeneous sensor networks operating in adversarial environments.

% --- Research Methods ---
\sect{Research Methods} The research methodology combines theoretical analysis, algorithm development, simulation, and hardware validation across four interconnected thrusts.

\subsect{Consensus Protocol Design} We will develop Byzantine fault-tolerant consensus protocols optimized for wireless sensor networks. Unlike traditional BFT protocols designed for data centers, our approach accounts for message loss, variable latency, and energy constraints inherent to embedded systems. The protocol will guarantee eventual consistency under the assumption that fewer than one-third of nodes exhibit Byzantine behavior.

Theoretical analysis will establish bounds on convergence time as a function of network diameter and message loss probability. We will prove safety and liveness properties using formal verification tools (TLA+ and Coq) to ensure correctness before implementation.

\subsect{Adaptive Topology Management} Sensor networks exhibit dynamic topology due to node failures, mobility, and environmental interference. We will design self-organizing protocols that maintain connectivity and load balance without centralized coordination. The approach leverages local neighbor discovery and gradient-based routing to establish hierarchical communication structures that adapt to changing conditions.

Machine learning techniques will predict topology changes based on historical patterns, enabling proactive route establishment before failures occur. Reinforcement learning agents at each node will learn optimal neighbor selection policies that balance communication reliability against energy consumption.

\subsect{Distributed Resource Allocation} Computational and communication resources in sensor networks are severely constrained. We will formulate resource allocation as a distributed optimization problem and develop algorithms based on alternating direction method of multipliers (ADMM) that converge to globally optimal solutions through local message passing.

The framework will support multiple optimization objectives including minimizing total energy consumption, maximizing sensing coverage, and balancing computational load across nodes. Game-theoretic analysis will characterize equilibrium properties when nodes have competing objectives.

\subsect{Hardware Testbed Validation} Theoretical results will be validated on a 50-node testbed comprising Raspberry Pi devices with LoRa wireless modules. The testbed will be deployed across the engineering building to evaluate performance in realistic indoor environments with multipath propagation and interference.

Experiments will measure latency, throughput, energy consumption, and fault tolerance under controlled failure injection. Comparison against baseline centralized and flooding-based approaches will quantify the benefits of our decentralized framework.

\noindent\textbf{Table 1: Hardware Components and Specifications}\\[0.2em]
\noindent\begin{tabular}{@{}p{0.25\textwidth}p{0.30\textwidth}p{0.40\textwidth}@{}}
\hline
\textbf{Component} & \textbf{Model} & \textbf{Key Specifications} \\
\hline
Processing Unit & Raspberry Pi 4B & 4-core ARM Cortex-A72, 4GB RAM \\
Wireless Module & Semtech SX1276 & LoRa 915MHz, 20dBm TX power \\
Power Supply & LiPo 3.7V 6000mAh & 12+ hours continuous operation \\
GPS Module & u-blox NEO-6M & 2.5m accuracy, 5Hz update rate \\
Environmental Sensor & BME280 & Temp, humidity, pressure \\
\hline
\end{tabular}

% --- Project Goals and Expected Results ---
\sect{Project Goals and Expected Results} The project will deliver both theoretical contributions and practical artifacts that advance the state of autonomous multi-agent systems.

\subsect{Theoretical Contributions} We expect to publish 3-4 peer-reviewed papers establishing: (1) tight bounds on consensus convergence in lossy networks, (2) competitive ratio analysis for online topology adaptation, (3) convergence guarantees for distributed ADMM under asynchronous updates, and (4) game-theoretic characterization of multi-objective resource allocation.

\subsect{Software Deliverables} Open-source implementations of all algorithms will be released under MIT license. The software package will include: consensus protocol library (Rust), topology management daemon (C++), resource allocation solver (Python), and simulation framework (Python/SimPy). Documentation will enable researchers to reproduce results and extend the framework.

\subsect{Hardware Prototype} The 50-node testbed will serve as a permanent research infrastructure available to the broader research community. Detailed construction guides and bill of materials will enable replication at other institutions.

\subsect{Quantitative Targets} Specific performance targets include: (1) consensus latency under 100ms for 100-node networks, (2) 99.9\% availability under 10\% simultaneous node failures, (3) 50\% reduction in energy consumption compared to flooding-based approaches, and (4) linear scaling of throughput with network size.

\noindent\textbf{Table 2: Performance Targets Comparison with Baseline Approaches}\\[0.2em]
\noindent\begin{tabular}{@{}p{0.35\textwidth}p{0.20\textwidth}p{0.20\textwidth}p{0.20\textwidth}@{}}
\hline
\textbf{Metric} & \textbf{Centralized} & \textbf{Flooding} & \textbf{Proposed} \\
\hline
Consensus Latency (ms) & 50 & 500 & $<$100 \\
Availability (\%) & 95.0 & 99.5 & 99.9 \\
Energy (mJ/message) & 12 & 85 & $<$40 \\
Scalability & O(n) & O($n^2$) & O(n log n) \\
Byzantine Tolerance & None & None & $<$1/3 \\
\hline
\end{tabular}

% --- Project Timeline ---
\sect{Project Timeline (12 Months)}

\vspace{-0.3em}
\begin{tabular}{@{}p{0.20\textwidth}p{0.75\textwidth}@{}}
\textbf{Months 1--3:} & Literature review, protocol specification, formal modeling \\[0.3em]
\textbf{Months 4--6:} & Algorithm implementation, simulation framework development \\[0.3em]
\textbf{Months 7--9:} & Hardware testbed construction, preliminary experiments \\[0.3em]
\textbf{Months 10--11:} & Full-scale validation, performance optimization \\[0.3em]
\textbf{Month 12:} & Documentation, paper submission, open-source release \\
\end{tabular}

% --- NASA Relevance ---
\sect{NASA Relevance} This research directly supports NASA's objectives in autonomous systems, distributed spacecraft, and planetary sensor networks. The consensus protocols developed here apply to formation flying missions where multiple spacecraft must coordinate without ground intervention. Topology adaptation algorithms address communication challenges in lunar and Martian surface networks where relay satellites provide intermittent connectivity.

The resource allocation framework enables autonomous prioritization of scientific observations when bandwidth to Earth is limited. As NASA missions increasingly rely on distributed architectures---from CubeSat constellations to heterogeneous rover teams---the theoretical foundations and practical tools developed in this project will accelerate mission development while improving reliability.

Collaboration with NASA Jet Propulsion Laboratory is planned to evaluate algorithm performance on representative deep-space communication scenarios. Results will inform requirements for future mission concept studies in distributed planetary science.

\end{document}
